% =====================================================================
\chapter{Méthodologie d'ingénierie}
\label{ch:methodo}
% =====================================================================

Un stage court impose une discipline : avec une durée limitée, la principale
cause d'échec n'est pas la difficulté technique, c'est la dispersion. Ce
chapitre décrit le cadre méthodologique mis en place, inspiré du développement
piloté par les spécifications (\emph{spec-driven development}), et les raisons
pour lesquelles il a été nécessaire.

\section{Le problème que la méthode résout}

Le projet a une particularité : une part importante du travail a été réalisée
en collaboration avec des agents de programmation assistée par IA. Ce mode de
travail est puissant pour la production de code, mais il introduit un risque
spécifique : l'agent qui reprend le travail ne connaît pas les décisions
antérieures, et peut \emph{défaire} un choix de conception sans s'en rendre
compte.

Deux exemples réels illustrent ce risque :

\begin{itemize}[leftmargin=1.4em]
  \item Un agent a réintroduit un niveau d'eau dépendant de la météo du jour.
        Conséquence : par temps sec, tous les scénarios affichaient 0~m, ce qui
        se lit comme une panne. Le comportement avait pourtant été
        explicitement corrigé.
  \item Un agent a utilisé une boîte englobante à l'échelle du quartier pour
        interroger le WFS. Mesuré à Nantes : la boîte se remplissait de
        polygones situés à 2~km, et le polygone contenant réellement l'adresse
        n'était jamais renvoyé dans la page de résultats. Résultat : \og hors
        TRI \fg{} sur un quai inondable.
\end{itemize}

Dans les deux cas, le code \emph{fonctionnait} au sens où il s'exécutait sans
erreur. Le défaut était sémantique. Une suite de tests ne suffit pas si les
tests eux-mêmes ne sont pas ancrés dans des décisions explicites.

\section{Le triptyque documentaire}

Trois documents structurent le projet, dans un ordre imposé.

\subsection{La constitution}

La \emph{constitution} (\fichier{constitution.md}) énonce les règles non
négociables. Chaque règle porte une étiquette d'application :

\begin{table}[H]
\centering
\small
\caption{Étiquettes d'application des règles constitutionnelles.}
\label{tab:enforcement}
\begin{tabularx}{\textwidth}{@{}l l X@{}}
\toprule
\textbf{Étiquette} & \textbf{Mécanisme} & \textbf{Exemple} \\
\midrule
\texttt{[CI]} & Vérification automatisée &
  Porte de recherche qui échoue si \code{niveau}, \code{score} ou
  \code{recommandations} apparaissent dans un schéma émis \\
\addlinespace
\texttt{[test]} & Assertion dans la suite &
  Le schéma refuse la sérialisation d'un objet portant un score, même si cet
  objet en contient un en interne \\
\addlinespace
\texttt{[review]} & Point de contrôle humain &
  Aucune nouvelle dépendance sans décision enregistrée dans le dépôt \\
\bottomrule
\end{tabularx}
\end{table}

La règle de révision est explicite : \emph{pour changer une règle, on amende la
constitution avant le changement qui la violerait}. Cela force la discussion à
se tenir au niveau du principe, pas au niveau du correctif.

\subsection{La spécification}

La spécification (\fichier{docs/workflow/spec.md}) décrit le produit en
exigences fonctionnelles numérotées, \frref{01} à \frref{28}, chacune assortie
de critères d'acceptation vérifiables. Par exemple :

\begin{description}[leftmargin=0em,style=nextline]
  \item[\frref{03}] Rejeter les géocodages de faible confiance avant tout appel
  aval : score $< 0{,}4 \rightarrow$ HTTP 422 \code{adresse\_ambigue} avec
  \code{label\_propose}. \emph{Critère :} une requête ambiguë renvoie 422 et
  aucun appel BDNB ou Géorisques n'a lieu (vérifié par compteurs de mocks).
  \item[\frref{14}] Ne jamais inférer une exposition depuis une source en
  échec : \code{present=true} uniquement depuis un test réussi,
  \code{false} depuis un test négatif, \code{null} si inconnu.
  \item[\frref{26}] Cache de résultats en mémoire, clé sur le triplet de
  géocodage normalisé (libellé + latitude/longitude arrondies à 5 décimales
  $\approx$ 1~m + code INSEE) et sur la version de schéma.
\end{description}

Les critères d'acceptation sont rédigés de manière à être directement
traduisibles en tests. C'est volontaire : une exigence sans test associé est un
souhait, pas une exigence.

\subsection{Le plan}

Le plan (\fichier{docs/workflow/plan.md}) découpe la spécification en tâches
T001 à T018, chacune avec ses dépendances, son estimation, ses points de
contact constitutionnels, ses critères d'acceptation et — point important —
\emph{les tests à écrire en premier}.

\begin{lstlisting}[caption={Extrait du plan : la tâche T008, porte de
certification de la v1.},label={lst:plan-t008}]
### T008 - Certification known-answer suite (v1 sign-off gate)
- Depends on: T004, T007
- Spec reference: FR-11
- Estimate: 4-8 h (fixture collection dominates)
- Acceptance criteria:
  - Recorded fixtures: known-inside and known-outside address per certified
    hazard (inondation/PPRI, mouvement_terrain, feu_foret, avalanche)
  - Inside -> present=true + resolution="per-building"; outside -> present=false
- Tests to write first:
  - test_certified_inside_<hazard> x4
  - test_certified_outside_<hazard> x4
\end{lstlisting}

\section{Le flux de travail}

Le flux impose un ordre qui n'est pas négociable :

\begin{center}
\begin{tikzpicture}[
  scale=0.76, transform shape,
  node distance=6mm,
  stepbox/.style={draw=typhoonBlue,fill=typhoonBlue!6,rounded corners=2pt,
               align=center,inner sep=5pt,font=\small,text width=2.5cm},
  arrow/.style={-{Latex[length=2mm]},thick,typhoonBlue}
]
\node[stepbox] (s1) {Brainstorm\\{\scriptsize décisions brutes}};
\node[stepbox,right=of s1] (s2) {Concept\\{\scriptsize questions résolues}};
\node[stepbox,right=of s2] (s3) {Spec\\{\scriptsize FR-01..FR-28}};
\node[stepbox,right=of s3] (s4) {Plan\\{\scriptsize T001..T018}};
\node[stepbox,right=of s4] (s5) {Test\\{\scriptsize commit qui échoue}};
\node[stepbox,right=of s5] (s6) {Code\\{\scriptsize commit qui passe}};
\foreach \a/\b in {s1/s2,s2/s3,s3/s4,s4/s5,s5/s6}
  \draw[arrow] (\a) -- (\b);
\end{tikzpicture}
\end{center}

Deux règles de ce flux méritent d'être soulignées.

\begin{enumerate}[leftmargin=1.6em]
  \item \textbf{Un changement de contrat passe d'abord par la constitution,
        puis par la spécification — jamais l'inverse.} Si une fonctionnalité
        nécessite de violer un invariant, ce n'est pas la fonctionnalité qu'on
        ajuste en premier, c'est le principe.
  \item \textbf{Les tests existants ne sont jamais modifiés pendant une tâche
        de fonctionnalité.} Un test jugé erroné est signalé et corrigé dans un
        change set séparé. Sans cette règle, il est toujours possible de faire
        passer un test en le réécrivant — ce qui vide la suite de son sens.
\end{enumerate}

\section{Convention de commit}

Les messages suivent une convention préfixée par le domaine :
\code{feat(etape2):}, \code{fix(flood):}, \code{test(SCN-xxx):},
\code{docs:}, \code{chore:}. Pour toute fonctionnalité, deux commits : un
commit \code{test(TXXX)} qui échoue, puis un commit \code{feat(TXXX)} qui
passe. Cette contrainte rend visible, dans l'historique, que le test a bien
précédé le code.

\section{Décisions documentées}

Les décisions qui poseraient problème à un lecteur futur font l'objet de
documents courts. Trois exemples du projet :

\begin{table}[H]
\centering
\small
\caption{Exemples de décisions documentées et leur motivation.}
\label{tab:adr}
\begin{tabularx}{\textwidth}{@{}l X@{}}
\toprule
\textbf{Décision} & \textbf{Document et motivation} \\
\midrule
Suppression du score composite &
  \fichier{concept.md}, OQ-1 : le score engage la responsabilité et se fait
  reconstruire ou refuser par l'actuaire \\
\addlinespace
RGA communal &
  \fichier{concept.md}, OQ-2 : la source vectorielle n'est pas câblée, annoncer
  plutôt que surinterpréter \\
\addlinespace
Licences commerciales &
  \fichier{concept.md}, OQ-3 : Géorisques et BDNB sont en Licence Ouverte 2.0,
  la réutilisation commerciale est permise, l'obligation d'attribution
  coïncide avec la conception de provenance \\
\addlinespace
Mode VFX &
  \fichier{constitution.md} §2.1, amendement du 14 septembre : autoriser un
  rendu cinématique à condition qu'il soit étiqueté non contractuel et
  strictement confiné à la présentation \\
\bottomrule
\end{tabularx}
\end{table}

\section{Outils}

\begin{table}[H]
\centering
\small
\caption{Outils et pratiques du projet.}
\label{tab:outils}
\begin{tabularx}{\textwidth}{@{}l X@{}}
\toprule
\textbf{Domaine} & \textbf{Choix} \\
\midrule
Gestion de version & Git, une tâche = une branche = une PR, messages préfixés \\
\addlinespace
Backend & Python 3.11+, FastAPI, Pydantic v2, \code{httpx} en asynchrone,
  \code{uvicorn} \\
\addlinespace
Frontend & TypeScript, React, Vite, Mapbox GL JS, Vitest \\
\addlinespace
Tests & \code{pytest} + \code{pytest-asyncio} (backend), Vitest (frontend) \\
\addlinespace
Lint / format & \code{ruff} (backend), \code{eslint} + \code{tsc -b} (frontend) \\
\addlinespace
Déploiement & Vercel (mono-projet : statique + fonction Python serverless) \\
\addlinespace
Documentation & Markdown dans le dépôt, lu par les humains et par les agents \\
\bottomrule
\end{tabularx}
\end{table}

\section{Apport de la méthode}

Le bénéfice concret de cette organisation a été de rendre les décisions
\emph{explicites et opposables}. Lorsqu'un comportement a été corrigé, la
correction est inscrite à la fois dans la constitution (le principe), dans la
spécification (l'exigence), dans le plan (la tâche) et dans la suite de tests
(l'assertion). Un agent — ou un humain — qui reprend le code trouve la règle
avant de trouver le piège.

C'est ce qui a permis de stabiliser un projet dont une part du code a été
écrite en collaboration avec des agents, sans que le produit ne dérive.
